% ═════════════════════════════════════════════════════════════════════════════
%  pages/resume.tex
% ═════════════════════════════════════════════════════════════════════════════
\chapter*{Résumé}
\addcontentsline{toc}{chapter}{Résumé}
\markboth{Résumé}{Résumé}

La multiplication des architectures microservices et le rythme croissant des cycles
de livraison logicielle soulèvent de nouveaux défis en matière de sécurité.
Le paradigme \textbf{DevSecOps} propose d'intégrer la sécurité à chaque étape
du cycle de développement, tandis que le \textbf{Service Mesh} offre une couche
d'infrastructure transparente pour sécuriser les communications interservices.

Ce rapport présente une étude approfondie de la combinaison de ces deux approches,
en mettant en lumière leur complémentarité pour garantir une sécurité avancée dans
les environnements distribués. Nous proposons notamment un prototype opérationnel
déployé sur \textbf{Kubernetes} avec \textbf{Istio}, intégrant le chiffrement
mutuel TLS (mTLS), l'authentification par jetons JWT et le contrôle d'accès
par rôles (RBAC), validé par des tests DAST et une analyse de menaces selon
la méthodologie STRIDE.

\medskip
\noindent\textbf{Mots-clés :} DevSecOps, Service Mesh, Istio, Kubernetes,
mTLS, DAST, SAST, microservices, sécurité, STRIDE, RBAC, JWT.

